\section{Teoremas de Fluxo}

\textbf{Notação.} $S$ é a fonte, $T$ é o sumidouro, $c(u, v)$ é a capacidade e $\text{maxflow}(S, T)$ é o fluxo
máximo. Depois do fluxo, $A$ é o que a BFS de $S$ alcança no residual (arestas com capacidade residual $> 0$).

\subsection*{Base}

\begin{itemize}
  \item \textbf{Corte mínimo:} $\text{maxflow}(S, T)$ é a menor soma de capacidades de arestas que, removidas,
        separam $S$ de $T$. O corte são as arestas originais $u \to v$ com $u \in A$ e $v \notin A$.
  \item \textbf{Aresta em todo corte mínimo:} $u \in A$ e $v$ alcança $T$ no residual.
  \item \textbf{Aresta em algum corte mínimo:} saturada e $v$ não é alcançável a partir de $u$ no residual.
  \item \textbf{Várias fontes ou sumidouros:} $S \to$ cada fonte e cada sumidouro $\to T$, com $\infty$.
  \item \textbf{Aresta não dirigida:} $u \to v$ e $v \to u$, ambas com $c$.
\end{itemize}

\subsection*{Reduções clássicas}

\begin{itemize}
  \item \textbf{Caminhos disjuntos em arestas:} capacidade $1$ em cada aresta. O fluxo é o maior número de caminhos
        de $S$ a $T$ sem aresta em comum, e também o menor número de arestas que separam $S$ de $T$.
  \item \textbf{Capacidade em vértice:} divida $u$ em $u_{in} \to u_{out}$ com capacidade $c(u)$; as arestas que entram
        vão para $u_{in}$ e as que saem partem de $u_{out}$. Com $c(u) = 1$, dá caminhos disjuntos em vértices. Com
        $c(u) = $ peso de $u$ e $\infty$ nas arestas, o corte é o conjunto de vértices de peso mínimo que separa $S$ de $T$.
  \item \textbf{Project selection:} escolher $u$ obriga escolher $v$. Lucro $p_u > 0$: $S \to u$ com $p_u$;
        custo $c_u > 0$: $u \to T$ com $c_u$; dependência: $u \to v$ com $\infty$.
        Lucro máximo $= \sum p_u - \text{corte mínimo}$, e os escolhidos são $A \setminus \{S\}$.
  \item \textbf{Escolha entre dois lados:} cada $u$ fica do lado de $S$ (custo $b_u$) ou de $T$ (custo $a_u$):
        $S \to u$ com $a_u$ e $u \to T$ com $b_u$. Penalidade $c$ se $u$ e $v$ ficam em lados diferentes:
        $u \to v$ e $v \to u$, ambas com $c$. O custo mínimo é o corte mínimo.
\end{itemize}

\subsection*{DAG com pesos nos vértices}

\begin{itemize}
  \item \textbf{Anticadeia de peso máximo:} no fecho transitivo, ligue $S \to u_L$ com $w_u$, $u_L \to v_R$ com
        $\infty$ se $u$ alcança $v$, e $v_R \to T$ com $w_v$. O peso máximo é $\sum w - \text{maxflow}(S, T)$, e a
        anticadeia são os $v$ com $v_L \in A$ e $v_R \notin A$.
  \item \textbf{Dilworth com pesos:} esse valor também é o menor número de caminhos (podendo repetir vértices) em que
        cada $v$ aparece pelo menos $w_v$ vezes. Para montar os caminhos, divida $v$ em $v_{in} \to v_{out}$ com
        limite inferior $w_v$, ligue $S \to v_{in}$, $v_{out} \to T$ e $u_{out} \to v_{in}$ para cada aresta, tudo com
        $\infty$, e calcule o fluxo mínimo.
\end{itemize}

\subsection*{Limite inferior nas arestas ($l \le f \le r$)}

Mande os $l$ obrigatórios em cada aresta e corrija o desequilíbrio que isso cria nos vértices.
\begin{itemize}
  \item Cada aresta $u \to v$ vira $u \to v$ com capacidade $r - l$, e o saldo muda: $d[v] \mathrel{+}= l$, $d[u] \mathrel{-}= l$.
  \item Adicione $T \to S$ com $\infty$, uma fonte nova $S'$ e um sumidouro novo $T'$: $S' \to v$ com $d[v]$ se $d[v] > 0$,
        e $v \to T'$ com $-d[v]$ se $d[v] < 0$.
  \item \textbf{Viável} $\iff$ $\text{maxflow}(S', T') = \sum_{d[v] > 0} d[v]$. O fluxo real de cada aresta é $l$ mais o fluxo
        nela na rede nova.
  \item \textbf{Fluxo máximo de $S$ a $T$:} seja $f_0$ o fluxo na aresta $T \to S$. Remova essa aresta e rode no residual:
        resposta $= f_0 + \text{maxflow}(S, T)$.
  \item \textbf{Fluxo mínimo:} igual, mas a resposta é $f_0 - \text{maxflow}(T, S)$.
\end{itemize}

\subsection*{Gomory-Hu}

Em grafo não dirigido, $n - 1$ fluxos montam uma árvore em que o corte mínimo entre $u$ e $v$ é a menor aresta do
caminho entre eles. Comece com $p[i] = 0$ para todo $i$. Para $i = 1, \dots, n - 1$: calcule o corte entre $i$ e
$p[i]$, com valor $c$ e lado $A$ de $i$; crie a aresta $(i, p[i], c)$; e, para todo $j > i$ com $p[j] = p[i]$ e
$j \in A$, faça $p[j] = i$.
